\documentclass[../main.tex]{subfiles}

\begin{document}

\title{\textbf{Harry Pyper Undergraduate application to:}\par From Simulation to AI: Physics-Informed Neural Networks for Chemical Kinetics in Turbulent Hydrogen Swirling Flames.}
\author{Umair Ahmed, School of Engineering}
\date{}
\maketitle
%\doublespacing
\setlength{\parskip}{1em}
\linespread{1.66}
\justifying
\section*{Degree course}

Mechanical Engineering BEng - Newcastle University, Tyne and Wear, UK

\section*{Interest in the subject}

Over the past two years, I have developed a strong interest in fluid mechanics and its applications in aerospace propulsion. In particular, I am fascinated by fluid mixing and the changes in fluid properties during combustion.\par

This interest has been shaped by early exposure to aviation, having grown up near Dunsfold Aerodrome and regularly attending air shows. The United Kingdom’s long-standing leadership in aerospace engineering has further strengthened my motivation to pursue this field.\par

My current university project focuses on computational fluid dynamics, where I perform Reynolds-Averaged Navier–Stokes (\textbf{RANS}) analysis of turbulent swirling jets using Python. This involves time-averaging and axisymmetric azimuthal averaging, integrating Fortran code, and working with \textbf{HDF5} data structures on Newcastle University’s SLURM-based \textbf{high-performance computing} system (COMET).\par

Through this work, I have developed strong computational and analytical skills and confirmed my ambition to pursue a career in computational fluid dynamics and aerospace propulsion.\par

\section*{Motivation for applying for an N8 CIR internship}

%Building communities of practice focused on N8 strategic research priorities;
Through my degree I have learned to appreciate the value in collaboration for solving difficult engineering problems. The N8 CIR shows a unique opportunity to work in a community that promotes cooperation for progress in research and the sharing of knowledge across disciplines.\par

This internship is particularly appealing as it combines computational fluid dynamics with emerging machine learning techniques. The opportunity to explore physics-informed neural networks in the context of turbulent combustion aligns directly with both my academic experience and future career goals.

%Developing CIR skills, including building a community of research software engineers; developing shared software infrastructure to support research priorities and cross-cutting capabilities;


%Facilitating the provision of shared high-performance computing resources involving N8 partners;

\section*{What I hope to gain from completing the internship}

I aim to deepen my understanding of machine learning methods, particularly physics-informed neural networks, and how they can be applied to complex fluid and combustion problems.\par

By contributing to this project, I hope to strengthen my computational and research skills while gaining experience in cutting-edge modelling techniques. This internship would represent an important step towards my long-term goal of working in aerospace propulsion engineering.\par


\end{document}
